\documentclass[10pt,a4paper]{amsart}
\usepackage[margin=19mm]{geometry}
\usepackage{amsmath,amssymb,mathtools}
\usepackage[scr=rsfs,cal=euler]{mathalfa}
\usepackage{xfrac}
\usepackage[unicode,hidelinks,bookmarks=false]{hyperref}
\newcommand{\ie}{i.e.\ }
\newcommand{\cf}{cf.\ }
\newcommand{\resp}{respectively\ }
\newcommand{\Subsection}{\subsection}
\providecommand{\nobreakdash}{\nobreak\hbox{-}}
\setlength{\emergencystretch}{2em}
\multlinegap0pt
\usepackage{tikz}
\usepackage{amssymb}
\usepackage{mathtools}
\newtheorem{proposition}{Proposition}
\newtheorem{remark}{Remark}
\usepackage{cleveref}
\crefname{proposition}{Proposition}{Propositions}
\crefname{remark}{Remark}{Remarks}
\newcommand{\Abs}[1]{\left| #1 \right|}
\newcommand{\Brack}[1]{\left( #1 \right)}
\newcommand{\Exp}[1]{\mathbb{E} #1}
\newcommand{\Prob}[1]{\mathbb{P}\left( #1 \right)}
\newcommand{\RR}{\mathbb{R}}
\newcommand{\Set}[1]{\left\{ #1 \right\}}
\newcommand{\bZ}{\bar{Z}}
\newcommand{\cond}{\mid}
\newcommand{\eps}{\varepsilon}
\newcommand{\mcN}{\mathcal{N}}
\newcommand{\mcX}{\mathcal{X}}
\newcommand{\norm}[1]{\left\|#1\right\|}
\title{How Deep Are Deep GPs, Really? A Sharp Threshold and a Non-Gaussian Limit for Compositional GPs}
\author{Mark Kozdoba, Shie Mannor}
\date{Source: arXiv:2606.08218v1 (CC BY 4.0)}
\begin{document}
\maketitle

\noindent\emph{Source and licence.} How Deep Are Deep GPs, Really? A Sharp Threshold and a Non-Gaussian Limit for Compositional GPs. Version arXiv:2606.08218v1 (CC BY 4.0), distributed by arXiv under the Creative Commons Attribution 4.0 International licence, http://creativecommons.org/licenses/by/4.0/, as recorded in the arXiv licence metadata for this exact version and shown on the abstract page https://arxiv.org/abs/2606.08218v1. Full licence notice: \texttt{licenses/arxiv-2606.08218-CC-BY-4.0.txt}.\par\smallskip

\noindent\textit{Editorial scope.} Counted unit: the article's proposition prop:sync\_v\_convergence (source0.tex lines 417--430). The article's complete original proof of it is delivered with it (source0.tex lines 432--551, one \texttt{proof} environment of 120 lines). No mathematics is written, repaired or re-derived: the statement and the proof are the article's own lines, in the article's own notation, and no retained line is substituted or re-flowed.  Retained uncounted as the source-local context the proof needs: lines 217--222; lines 410--415; lines 553--572.  Nothing was omitted from any retained span, so the package carries no omission record.  The retained lines cite 2 works, whose entries are transcribed verbatim from the article's own bibliography source in the e-print and delivered with short editorial labels recorded in the item's reference mapping.  No retained line contains a figure environment, an image-inclusion command or drawing code, so no image is delivered and the package declares no asset dependency.  Editorial formatting only, in addition: an AMS \texttt{amsart} preamble that restates the article's own declarations and macros used by the retained lines, namely the environments \texttt{proposition}, \texttt{remark} on separate counters, together with the standard packages those lines need.  The retained environments print in this document as Proposition 1, Remark 1 in order of appearance, whereas the article prints the same items with the numbering of its own full document.  The complete native source of the article as distributed in the arXiv e-print for this exact version is bundled unchanged as \texttt{sources/202317/source0.tex}.
The first object we study is the pairwise distance for a fixed pair of coordinates $s \ne t$, held fixed throughout this and the next subsection. Let $V_i := \bZ_{i,s} - \bZ_{i,t} \in \RR^d$ and $v_i := \norm{V_i}$. Conditional on $\bZ_i$, the pair $(\bZ_{i+1,s}, \bZ_{i+1,t})$ is jointly centred Gaussian with cross-covariance $\phi(v_i)\, I_d$, so
\begin{eqnarray}
  \label{eq:kernel_corr_converegnce_d}
  v_{i+1}^2 \;=\; 2\Brack{1 - \phi(v_i)}\, X_i, \qquad X_i \sim \chi^2_d,
\end{eqnarray}
with $X_i$ independent of $v_i$ by rotation invariance. Because the $G_i$ are i.i.d.\ across $i$, the innovations $\Set{X_i}_{i \ge 1}$ at this fixed pair are i.i.d., and the scalar chain $\Set{v_i^2}$ is Markov on $[0, \infty)$.

For $d = 1$, the recursion \eqref{eq:kernel_corr_converegnce_d} reduces to
\begin{eqnarray}
  \label{eq:kernel_corr_converegnce}
  v_{i+1} \;=\; \sqrt{2(1 - \phi(v_i))}\, g_i, \qquad g_i \sim \mcN(0, 1) \text{ i.i.d.},
\end{eqnarray}
with $v_i := Z_{i,1} - Z_{i,2} \in \RR$.

\begin{remark}
  \label{rem:log_g_sq_integrable}
  For $g \sim \mcN(0, 1)$, $\Exp{|\log g^2|} < \infty$ (used in the proof of (i) and in the non-convergence step of (ii)). Indeed, $g^2 \sim \chi^2_1$ has density $f(x) = (2\pi x)^{-1/2} e^{-x/2}$ on $(0, \infty)$, and
  \begin{eqnarray*}
    \Exp{|\log g^2|} \;=\; \int_0^\infty |\log x|\,(2\pi x)^{-1/2} e^{-x/2}\,dx.
  \end{eqnarray*}
  Split at $x = 1$. The upper tail $\int_1^\infty (\log x)(2\pi x)^{-1/2} e^{-x/2}\,dx$ is finite because $(\log x)\,x^{-1/2} e^{-x/2}$ decays exponentially. For the lower tail, $|\log x| = -\log x$ on $(0,1)$ and $e^{-x/2} \le 1$, so
  \begin{eqnarray*}
    \int_0^1 (-\log x)\,(2\pi x)^{-1/2} e^{-x/2}\,dx \;\le\; (2\pi)^{-1/2}\!\int_0^1 \frac{-\log x}{\sqrt{x}}\,dx.
  \end{eqnarray*}
  Substitute $u := -\log x$, so $x = e^{-u}$, $dx = -e^{-u}\,du$, $\sqrt{x} = e^{-u/2}$, and the limits $x \in (0,1)$ map to $u \in (\infty, 0)$. The integrand transforms as
  \begin{eqnarray*}
    \frac{-\log x}{\sqrt{x}}\,dx \;=\; \frac{u}{e^{-u/2}}\cdot(-e^{-u}\,du) \;=\; -u\,e^{-u/2}\,du,
  \end{eqnarray*}
  and the sign flip cancels against the reversed limits, yielding
  \begin{eqnarray*}
    \int_0^1\frac{-\log x}{\sqrt{x}}\,dx \;=\; \int_0^\infty u\,e^{-u/2}\,du \;=\; 4.
  \end{eqnarray*}
  The substitution thus replaces the logarithmic singularity at $x = 0$ by a linear factor $u$ against an exponentially decaying weight, which is harmless. Combined with finiteness of the upper tail, this gives $\Exp{|\log g^2|} < \infty$, and a fortiori $\Exp{\log g^2}$ is well-defined and equals $-\gamma - \log 2$.
\end{remark}

\begin{proposition}[Scalar dichotomy, $d = 1$]
  \label{prop:sync_v_convergence}
  Let $\Set{v_i}_{i \ge 1}$ be the chain defined by \eqref{eq:kernel_corr_converegnce} with $v_1 \in \RR$ arbitrary, and let
  \begin{eqnarray*}
    r_c \;:=\; \frac{1}{\sqrt{2}\, e^{\gamma/2}} \;=\; \sqrt{\frac{1}{2 e^\gamma}} \;\approx\; 0.5298.
  \end{eqnarray*}
  \begin{itemize}
    \item[(i)] If $r > r_c$, then for every $v_1 \in \RR$, $v_i \to 0$ almost surely, and
      \begin{eqnarray*}
        \limsup_{i \to \infty} \frac{1}{i}\, \log v_i^2 \;\le\; -\gamma - \log 2 - 2\log r \;<\; 0 \quad \text{a.s.}
      \end{eqnarray*}
    \item[(ii)] If $r < r_c$ and $v_1 \ne 0$, then $v_i \not\to 0$ almost surely, and $\Set{v_i^2}$ admits a unique nontrivial stationary distribution on $(0, \infty)$, independent of $v_1$.
  \end{itemize}
\end{proposition}

\begin{proof}[Proof of \Cref{prop:sync_v_convergence}]
  Set $u_i := v_i^2 \ge 0$, $L_i := \log u_i$, and $F(u) := 2(1 - e^{-u/(2r^2)})$. Then \eqref{eq:kernel_corr_converegnce} reads
  \begin{eqnarray}
    \label{eq:u_chain}
    u_{i+1} \;=\; F(u_i)\,g_i^2.
  \end{eqnarray}
  If $u_1 = 0$ then $u_i \equiv 0$; otherwise $u_i > 0$ a.s.\ for all $i$ (since $g_i \ne 0$ a.s.), and \eqref{eq:u_chain} gives
  \begin{eqnarray}
    \label{eq:L_chain}
    L_{i+1} & = & L_i + H(L_i) + \log g_i^2,
  \end{eqnarray}
  where $H(L) := \log\Brack{F(e^L)/e^L}$.
  The function $H$ is continuous on $\RR$, satisfies $H(L) \to -2\log r$ as $L \to -\infty$ (since $F(u)/u \to 1/r^2$ as $u \to 0$), and obeys the uniform upper bound
  \begin{eqnarray}
    \label{eq:H_upper}
    H(L) \;\le\; -2\log r \qquad \text{for all } L \in \RR,
  \end{eqnarray}
  which follows from $1 - e^{-x} \le x$. For $g \sim \mcN(0,1)$, $\Exp{\log g^2} = \psi(1/2) + \log 2 = -\gamma - \log 2$; set $\rho := -\gamma - \log 2 - 2\log r$, so that $r > r_c$ iff $\rho < 0$, and $r < r_c$ iff $\rho > 0$.

  \emph{Proof of (i).} Assume $r > r_c$, so $\rho < 0$. From \eqref{eq:L_chain} and \eqref{eq:H_upper},
  \begin{eqnarray*}
    L_i \;\le\; L_1 + \sum_{j=1}^{i-1}\Brack{-2\log r + \log g_j^2}.
  \end{eqnarray*}
  We invoke Kolmogorov's strong law of large numbers for i.i.d.\ sequences (Theorem~2.4.1 of \cite{durrett2019probability}): if $\Set{Y_j}_{j \ge 1}$ are i.i.d.\ with $\Exp{|Y_1|} < \infty$, then $n^{-1}\sum_{j=1}^n Y_j \to \Exp{Y_1}$ almost surely. With $Y_j := -2\log r + \log g_j^2$ the sequence is i.i.d.\ (since the $g_j$ are), and the integrability hypothesis $\Exp{|\log g^2|} < \infty$ holds by \Cref{rem:log_g_sq_integrable}. The SLLN therefore yields $(i-1)^{-1}\sum_{j=1}^{i-1}(-2\log r + \log g_j^2) \to \rho$ a.s., whence $\limsup_{i\to\infty} L_i/i \le \rho < 0$ a.s. So $u_i \to 0$ a.s.\ at exponential rate at least $|\rho|$, i.e.\ $v_i \to 0$ a.s.

  \emph{Proof of (ii), non-convergence.} Assume $r < r_c$, so $\rho > 0$. Fix $\eps \in (0, \rho/2)$. Since $H(L) \to -2\log r$ as $L \to -\infty$, there exists $L_* < 0$ such that
  \begin{eqnarray}
    \label{eq:H_lower}
    H(L) \;\ge\; -2\log r - \eps \qquad \text{for all } L \le L_*;
  \end{eqnarray}
  write $\delta := e^{L_*} > 0$. Fix $u_1 > 0$ and define the hitting time
  \begin{eqnarray*}
    \tau \;:=\; \inf\Set{i \ge 1 : u_i > \delta}.
  \end{eqnarray*}
  We claim $\tau < \infty$ a.s. On the event $\Set{\tau = \infty}$, $L_i \le L_*$ for every $i$, so \eqref{eq:L_chain} and \eqref{eq:H_lower} give the pathwise lower bound
  \begin{eqnarray}
    \label{eq:L_lower_on_tau}
    L_{n+1} & \ge & L_1 + \sum_{j=1}^n\Brack{-2\log r - \eps + \log g_j^2} \nonumber \\
    & = & L_1 + n\Brack{\rho - \eps} + S_n,
  \end{eqnarray}
  where $S_n := \sum_{j=1}^n(\log g_j^2 - \Exp{\log g^2})$ is a zero-mean random walk with i.i.d.\ increments. By the same SLLN as in part~(i), applied \emph{unconditionally} to the i.i.d.\ sequence $\Set{\log g_j^2 - \Exp{\log g^2}}_{j \ge 1}$ on the underlying probability space, $\Prob{\Omega_0} = 1$ for $\Omega_0 := \Set{S_n/n \to 0}$; we do \emph{not} condition on $\Set{\tau = \infty}$ before invoking the SLLN, since under that conditioning the $g_j$ would no longer be i.i.d. Now consider any sample point $\omega \in \Set{\tau = \infty} \cap \Omega_0$. The pathwise bound \eqref{eq:L_lower_on_tau} together with $S_n(\omega)/n \to 0$ and $\rho - \eps > \rho/2 > 0$ force $L_{n+1}(\omega) \to +\infty$, contradicting $L_{n+1}(\omega) \le L_*$ for all $n$ (which holds since $\omega \in \Set{\tau = \infty}$). Hence $\Set{\tau = \infty} \cap \Omega_0 = \emptyset$, and combined with $\Prob{\Omega_0} = 1$ this gives $\Prob{\tau = \infty} = 0$.

  Iterate: set $\tau_0 := 0$ and $\tau_{k+1} := \inf\Set{i > \tau_k : u_i > \delta}$. By the strong Markov property applied at $\tau_k + 1$ (from which point the chain restarts from the positive state $u_{\tau_k + 1}$), the same argument gives $\tau_{k+1} < \infty$ a.s.\ on $\Set{\tau_k < \infty}$. By induction every $\tau_k$ is a.s.\ finite, so $u_i > \delta$ for infinitely many $i$ a.s.; in particular
  \begin{eqnarray*}
    \limsup_{i\to\infty} u_i \;\ge\; \delta \;>\; 0 \qquad \text{a.s.,}
  \end{eqnarray*}
  so $v_i = \pm\sqrt{u_i}$ does \emph{not} converge to $0$ a.s.

  \emph{Proof of (ii), stationary distribution.} We establish positive Harris recurrence via a Foster--Lyapunov drift argument; we briefly recall the framework before applying it.

  Let $\Set{X_i}$ be a Markov chain on a state space $\mcX$ with transition kernel $P(x, \cdot)$. The chain is \emph{$\psi$-irreducible} if there exists a non-trivial measure $\psi$ on $\mcX$ such that for every $x \in \mcX$ and every measurable $A$ with $\psi(A) > 0$, the chain reaches $A$ from $x$ with positive probability in some finite number of steps; intuitively, no part of the state space is invisible from any starting point. A set $C \subset \mcX$ is \emph{small} if there exist $m \ge 1$, $\eta > 0$, and a probability measure $\nu$ such that $P^m(x, \cdot) \ge \eta\,\nu(\cdot)$ for every $x \in C$, i.e.\ $C$ is a uniform regeneration region: from anywhere in $C$, the $m$-step law dominates a fixed probability measure. For a $\psi$-irreducible chain on a general state space there is a notion of \emph{period} $d \ge 1$, defined via the $d$-cycle of small sets associated with the chain (see Section~5.4 of \cite{meyn1993markov}, where the period is defined for general state-space $\psi$-irreducible chains), and the chain is \emph{aperiodic} precisely when $d = 1$. A convenient sufficient condition --- which we will use --- is that some small set $C$ admit a one-step minorisation $P(x, \cdot) \ge \eta\,\nu$ for every $x \in C$ (i.e.\ the minorising step $m$ above can be taken equal to $1$); a chain satisfying this is called \emph{strongly aperiodic}, and strong aperiodicity implies aperiodicity.

  Two results from \cite{meyn1993markov} combine to give positive recurrence and convergence to equilibrium under a Lyapunov drift condition. First, \emph{Foster's drift criterion} (Theorem~11.3.4 of \cite{meyn1993markov}): if the chain is $\psi$-irreducible and there exist a measurable function $V : \mcX \to [0, \infty)$, a small set $C$, and constants $c > 0$, $b < \infty$, such that the drift condition
  \begin{eqnarray}
    \label{eq:foster_lyapunov_drift}
    \Exp{V(X_{i+1}) \cond X_i = x} \;\le\; V(x) - c + b\,\mathbf{1}_C(x)
  \end{eqnarray}
  holds for every $x \in \mcX$, then the chain is positive Harris recurrent and admits a unique invariant probability measure $\pi$. Second, the \emph{Aperiodic Ergodic Theorem} (Theorem~13.0.1 of \cite{meyn1993markov}): once the chain is positive Harris recurrent, aperiodic, and has invariant probability $\pi$, the law of $X_i$ converges to $\pi$ in total variation from every starting point $x \in \mcX$,
  \begin{eqnarray*}
    \norm{P^i(x, \cdot) - \pi}_{TV} \to 0 \qquad \text{as } i \to \infty.
  \end{eqnarray*}
  The function $V$ plays the role of a Lyapunov function: it measures a generalised ``distance'' to the small set $C$, and \eqref{eq:foster_lyapunov_drift} asserts that this distance shrinks on average by at least $c$ outside $C$, while inside $C$ it may grow but only by a bounded amount $b$. Together these prevent both escape to infinity and trapping at the boundary of $\mcX$, forcing the chain to return to $C$ infinitely often with controlled hitting times --- which is exactly what positive recurrence requires.

  We apply this framework to the chain $\Set{u_i}$ on $\mcX = (0, \infty)$ with transition $u \mapsto F(u)\,g^2$. Since $g^2$ has a smooth strictly positive Lebesgue density on $(0,\infty)$ and $F(u) > 0$ for $u > 0$, the one-step transition law of $u_{i+1}$ given $u_i = u$ has a strictly positive Lebesgue density on $(0,\infty)$, smoothly depending on $u$. This has three consequences. \emph{(a)} The chain is $\psi$-irreducible with respect to Lebesgue measure on $(0, \infty)$: any set of positive Lebesgue measure has positive one-step transition probability from every $u \in (0,\infty)$. \emph{(b)} Every compact $K \subset (0,\infty)$ is a small set with $m = 1$: the one-step density is bounded below by some $\eta_K > 0$ uniformly on $K \times K$ by continuity, hence $P(u, \cdot) \ge \eta_K\,|K|\cdot \mathrm{Unif}(K)$ for every $u \in K$. \emph{(c)} The chain is aperiodic: $m = 1$ in (b) gives the strong aperiodicity condition. Since $L = \log u$ is a homeomorphism between $(0,\infty)$ and $\RR$, all of these properties transfer verbatim to the log-chain $\Set{L_i}$ on $\RR$ (with Lebesgue measure pulled back to Lebesgue on $\RR$).

  For the Lyapunov function we work in log-coordinates, where the multiplicative dynamics become additive (cf.\ \eqref{eq:L_chain}). Define
  \begin{eqnarray*}
    V(L) \;:=\; (L_0 - L)^+ + (L - L^0)^+,
  \end{eqnarray*}
  where $L_0 > L_*$ and $L^0 > L_0$ are to be fixed. This is a ``tent'' function: $V \equiv 0$ on the strip $L \in [L_0, L^0]$, and $V$ grows linearly outside it (with slope $-1$ to the left of $L_0$ and slope $+1$ to the right of $L^0$). Geometrically, $V(L)$ is the distance from $L$ to the strip $[L_0, L^0]$, which in $u$-coordinates is the compact interval $[e^{L_0}, e^{L^0}] \subset (0,\infty)$ --- a small set by the previous paragraph. Verifying \eqref{eq:foster_lyapunov_drift} thus amounts to showing that $\Exp{V(L_{i+1}) \cond L_i = L} < V(L) - c$ for $L$ in the two tails $\Set{L \le L_*}$ and $\Set{L \ge L^0 + 2\kappa}$, with the in-between region $C := [L_*, L^0 + 2\kappa]$ acting as the small set on which $V$ is allowed to grow by the constant $b$. The mechanisms in the two tails are different: on the \emph{left} ($L \le L_*$) the linearisation-based positive drift $\rho - \eps$ on $L$ from \eqref{eq:H_lower} pushes $L_{i+1}$ rightward, decreasing $V$; on the \emph{right} ($L \ge L^0$) the global cap $u_{i+1} \le 2 g_i^2$ resets $L_{i+1}$ to a state-independent random level of bounded mean, which sits far below $L^0$ once $L^0$ is large, again decreasing $V$. We verify these now.

  \emph{Region 1: $L \le L_*$.} Here $V(L) = L_0 - L$. Using $(a)^+ = a + (-a)^+$,
  \begin{eqnarray*}
    \lefteqn{\Exp{(L_0 - L_{i+1})^+ \cond L_i = L}} \\
    & = & L_0 - \Exp{L_{i+1} \cond L_i = L} \\
    & & {}+ \Exp{(L_{i+1} - L_0)^+ \cond L_i = L}.
  \end{eqnarray*}
  By \eqref{eq:L_chain} and \eqref{eq:H_lower}, $\Exp{L_{i+1} \cond L_i = L} = L + H(L) - \gamma - \log 2 \ge L + (\rho - \eps)$. By \eqref{eq:H_upper}, $L_{i+1} \le L - 2\log r + \log g_i^2$, so for $L \le L_*$,
  \begin{eqnarray*}
    \Exp{(L_{i+1} - L_0)^+ \cond L_i = L} & \le & \Exp{(\log g^2 - M)^+},
  \end{eqnarray*}
  where $M := L_0 - L_* + 2\log r$.
  Since $\log g^2$ has finite mean, $\Exp{(\log g^2 - M)^+} \to 0$ as $M \to \infty$; choose $L_0$ large enough that this is $\le (\rho - \eps)/4$, and by the same bound applied with $L^0$ in place of $L_0$, also $\Exp{(L_{i+1} - L^0)^+ \cond L_i = L} \le (\rho - \eps)/4$ for $L \le L_*$. Using the identity $(L_0 - L_{i+1})^+ = (L_0 - L_{i+1}) + (L_{i+1} - L_0)^+$ together with $L_0 - \Exp{L_{i+1} \cond L_i = L} \le L_0 - L - (\rho - \eps) = V(L) - (\rho - \eps)$ (valid since $V(L) = L_0 - L$ for $L \le L_*$), we obtain
  \begin{eqnarray*}
    \lefteqn{\Exp{(L_0 - L_{i+1})^+ \cond L_i = L}} \\
    & \le & V(L) - (\rho - \eps) + (\rho - \eps)/4 \\
    & = & V(L) - 3(\rho - \eps)/4.
  \end{eqnarray*}
  Adding $\Exp{(L_{i+1} - L^0)^+ \cond L_i = L} \le (\rho - \eps)/4$ gives the drift bound
  \begin{eqnarray}
    \label{eq:drift_left}
    \Exp{V(L_{i+1}) \cond L_i = L} & \le & V(L) - (\rho - \eps)/2
  \end{eqnarray}
  for all $L \le L_*$, i.e.\ $c_1 := (\rho - \eps)/2 > 0$ on $\Set{L \le L_*}$.

  \emph{Region 2: $L \ge L^0$.} Here $V(L) = L - L^0$. From \eqref{eq:u_chain}, $u_{i+1} \le 2 g_i^2$, so $L_{i+1} \le \log 2 + \log g_i^2$ regardless of $L_i$. On Region 2 the lower bound $u_{i+1} \ge F(e^{L^0})\, g_i^2$ also holds (since $F$ is increasing and $u_i \ge e^{L^0}$), giving $L_{i+1} \ge \log F(e^{L^0}) + \log g_i^2$. Combining, for every $L_i \ge L^0$ the new state satisfies
  \begin{eqnarray*}
    \Abs{L_{i+1}} \;\le\; \max\!\Brack{\log 2,\, -\log F(e^{L^0})} + \Abs{\log g_i^2}
    \;=:\; W_i,
  \end{eqnarray*}
  an envelope independent of $L_i$ with finite mean (since $\Exp{|\log g^2|} < \infty$ by \Cref{rem:log_g_sq_integrable}). Hence
  \begin{eqnarray*}
    V(L_{i+1}) \;\le\; \Abs{L_0 - L_{i+1}} + \Abs{L_{i+1} - L^0} \;\le\; 2 W_i + L_0 + L^0,
  \end{eqnarray*}
  so that
  \begin{eqnarray*}
    \kappa \;:=\; \Exp{2 W_i + L_0 + L^0} \;<\; \infty,
  \end{eqnarray*}
  and $\Exp{V(L_{i+1}) \cond L_i = L} \le \kappa$ for every $L \ge L^0$. The constant $\kappa$ is fixed by the chain and the choices of $L_0, L^0$.

  Now we use this $\kappa$ to delimit the small set. We want a strictly negative drift $\le -\kappa$ in Region 2, i.e.\ $\Exp{V(L_{i+1}) \cond L_i = L} \le V(L) - \kappa$, which (using the uniform bound above) holds whenever $V(L) \ge 2\kappa$. Since $V(L) = L - L^0$ in Region 2, this requires $L \ge L^0 + 2\kappa$ --- and this is precisely why the small set is taken to extend up to $L^0 + 2\kappa$ rather than just to $L^0$: the strip $[L^0, L^0 + 2\kappa]$ is a buffer in which $V(L)$ is too small (between $0$ and $2\kappa$) for the uniform-envelope bound to dominate it, so we cannot yet claim a useful negative drift there. Outside this buffer, on $\Set{L \ge L^0 + 2\kappa}$, we have the drift $\Exp{V(L_{i+1}) \cond L_i = L} \le \kappa \le V(L) - \kappa$.

  \emph{Region 3: the compact set $C := [L_*, L^0 + 2\kappa]$.} Here $V(L) \le \max(L_0 - L_*,\, 2\kappa)$. The map $L \mapsto \Exp{V(L_{i+1}) \cond L_i = L}$ is continuous in $L$ (the transition density is continuous and has finite first moment of $|L_{i+1}|$ locally in $L$, since $L_{i+1} = \log F(e^L) + \log g^2$ with $\log F(e^L)$ continuous and $\Exp{|\log g^2|} < \infty$), hence bounded on the compact $C$ by some $b < \infty$. On $C$ we make no claim of a negative drift; this is the small set on which $V$ is allowed to grow, contributing the term $b\,\mathbf{1}_C(L)$ in \eqref{eq:foster_lyapunov_drift}.

  Combining the three regions yields the drift condition
  \begin{eqnarray*}
    \Exp{V(L_{i+1}) \cond L_i = L} \;\le\; V(L) - c + b\,\mathbf{1}_C(L),
  \end{eqnarray*}
  for some $c > 0$ and $b < \infty$, exactly the form \eqref{eq:foster_lyapunov_drift}, with small set $C = [L_*, L^0 + 2\kappa]$. Together with the $\psi$-irreducibility verified above, Foster's drift criterion (Theorem~11.3.4 of \cite{meyn1993markov}) yields that $\Set{u_i}$ is positive Harris recurrent and admits a unique invariant probability measure $\pi$ on $(0, \infty)$. Aperiodicity was also verified above ($m = 1$ minorisation on every compact). The Aperiodic Ergodic Theorem (Theorem~13.0.1 of \cite{meyn1993markov}) then applies, giving that the law of $u_i$ converges to $\pi$ in total variation from every $u_1 > 0$; in particular $\pi$ is independent of $v_1$. Non-degeneracy $\pi((\delta,\infty)) > 0$ follows from the non-convergence step (the chain spends a positive fraction of time above $\delta$).
\end{proof}

\begin{thebibliography}{99}
\bibitem[durret]{durrett2019probability} Rick Durrett. \newblock \emph{Probability: Theory and Examples}. \newblock Cambridge Series in Statistical and Probabilistic Mathematics. Cambridge University Press, 5 edition, 2019.
\bibitem[meyn19]{meyn1993markov} Sean~P. Meyn and Richard~L. Tweedie. \newblock \emph{Markov Chains and Stochastic Stability}. \newblock Communications and Control Engineering Series. Springer-Verlag London, 1993.
\end{thebibliography}
\end{document}
